\unitopen{04}
\begin{candos}
\item powiedzieć, co to jest (\de{Das ist \dots}), używając \de{der/die/das} i \de{ein/eine/ein};
\item odgadnąć rodzaj większości rzeczowników po końcówce;
\item utworzyć liczbę mnogą większości rzeczowników z tej jednostki;
\item opisać biurko i stanowisko pracy po niemiecku.
\end{candos}

\section{Słownictwo}
\vocabtable{04}

\section{Trzy rodzaje, jeden rzeczownik wielką literą}
Niemiecki ma \textbf{trzy rodzaje}: \artm{der} (męski), \artf{die} (żeński), \artn{das} (nijaki), a liczba mnoga
ma zawsze \artf{die}. \textbf{Każdy rzeczownik piszesz wielką literą} i uczysz się go z rodzajnikiem,
bo rodzaj nie zawsze zgadza się z polskim: \de{der Löffel} (łyżka --- polski: żeński), \de{das Mädchen}
(dziewczynka --- nijaki), \de{die Sonne} (słońce --- polski: nijaki).

\begin{center}\small
\begin{tabular}{@{}lll@{}}
\toprule
 & \textsf{\textbf{Określony}} & \textsf{\textbf{Nieokreślony}}\\
\midrule
męski & \artm{der} \de{Tisch} & \de{ein Tisch}\\
żeński & \artf{die} \de{Tasche} & \de{eine Tasche}\\
nijaki & \artn{das} \de{Buch} & \de{ein Buch}\\
mnoga & \artf{die} \de{Bücher} & (brak): \de{Bücher}\\
\bottomrule
\end{tabular}
\end{center}

\de{Das ist ein Computer.} (pierwsza wzmianka) \de{Der Computer ist neu.} (już wiesz, który).

\section{Rodzaj po końcówce}
\begin{center}\small
\begin{tabular}{@{}p{0.27\linewidth}p{0.17\linewidth}p{0.45\linewidth}@{}}
\toprule
\textsf{\textbf{Końcówka}} & \textsf{\textbf{Rodzaj}} & \textsf{\textbf{Przykład}}\\
\midrule
\de{-ung}, \de{-heit}, \de{-keit}, \de{-schaft}, \de{-ion}, \de{-tät}, \de{-ur} & \artf{die} & \de{die Anwendung}, \de{die Tastatur}, \de{die Funktion}\\
\de{-er} (wykonawca, urządzenie), \de{-ig}, \de{-ling} & \artm{der} & \de{der Computer}, \de{der Server}, \de{der Drucker}\\
\de{-chen}, \de{-lein}, \de{-ment}, \de{-um} & \artn{das} & \de{das Mädchen}, \de{das Dokument}, \de{das Zentrum}\\
dni, miesiące, pory roku & \artm{der} & \de{der Montag}, \de{der Januar}, \de{der Sommer}\\
\bottomrule
\end{tabular}
\end{center}
To reguły \emph{ze statystyki}, nie prawa: \de{der Fehler} i \de{das Passwort} łamią regułę końcówki.
Rodzaj złożenia wyznacza ostatni człon: \de{die Bank} $\rightarrow$ \de{die Datenbank}; \de{das Werk}
$\rightarrow$ \de{das Netzwerk}.

\section{Liczba mnoga: siedem wzorców}
\begin{center}\small
\begin{tabular}{@{}p{0.22\linewidth}p{0.36\linewidth}p{0.33\linewidth}@{}}
\toprule
\textsf{\textbf{Wzorzec}} & \textsf{\textbf{Przykład}} & \textsf{\textbf{Uwaga}}\\
\midrule
\de{-e} & \de{der Tisch $\rightarrow$ die Tische} & często dla rodzaju męskiego\\
\de{¨-e} & \de{der Stuhl $\rightarrow$ die Stühle} & przegłos: a, o, u $\rightarrow$ ä, ö, ü\\
\de{-er}, \de{¨-er} & \de{das Buch $\rightarrow$ die Bücher} & tylko rodzaj nijaki i kilka męskich\\
\de{-n}, \de{-en} & \de{die Tasche $\rightarrow$ die Taschen} & prawie wszystkie żeńskie\\
\de{-s} & \de{das Auto $\rightarrow$ die Autos} & zapożyczenia\\
\de{-} (bez zmiany) & \de{der Server $\rightarrow$ die Server} & rzeczowniki na \de{-er}, \de{-el}\\
\de{¨-} & \de{der Bruder $\rightarrow$ die Brüder} & przegłos bez końcówki\\
\bottomrule
\end{tabular}
\end{center}
Nie zgadujesz liczby mnogiej: uczysz się jej \emph{razem} z rzeczownikiem. Dlatego w słowniku jest
\de{der Tisch, Tische}.

\begin{falle}
Polskie rodzajniki nie istnieją, więc nie ma nic do „przetłumaczenia”. Nie szukaj logiki w \de{das Mädchen}
(płeć naturalna nie wygrywa z końcówką \de{-chen}). Drugi błąd to pisanie rzeczowników małą literą w e-mailu
(\de{ich habe einen fehler}): w niemieckim to błąd ortograficzny, nie styl.
\end{falle}

\begin{praca}
\de{Der Server läuft nicht.} \de{Die Datenbank ist langsam.} \de{Das Netzwerk ist ausgefallen.}
\de{Die Software hat einen Fehler.} Trzy zdania, które usłyszysz w pierwszym tygodniu pracy.
Kod i nazwy techniczne (\de{Code}, \de{Branch}) zostają po angielsku, ale rodzajnik jest niemiecki:
\de{der Code}, \de{der Branch}, \de{die Pipeline}, \de{das Repository}.
\end{praca}

\section{Dialogi}
\textbf{Dialog 1 --- w biurze, pierwszy dzień}
\begin{dialog}
\turn{A}{Das ist Ihr Arbeitsplatz. Hier ist der Computer und da ist der Bildschirm.}{To jest Pana stanowisko pracy. Tu jest komputer, a tam ekran.}
\turn{B}{Und was ist das?}{A co to jest?}
\turn{A}{Das ist eine Tastatur. Das ist ein Kabel.}{To jest klawiatura. To jest kabel.}
\turn{B}{Wo ist die Maus?}{Gdzie jest mysz?}
\turn{A}{Die Maus ist hier. Die Maus und die Tastatur sind neu.}{Mysz jest tutaj. Mysz i klawiatura są nowe.}
\end{dialog}

\section{Ćwiczenia}
\begin{cwiczenia}
\exitem{Wstaw rodzajnik określony: \de{\dots{} Tisch}.}{\artm{der} \de{Tisch}}
\exitem{Wstaw rodzajnik określony: \de{\dots{} Datei}.}{\artf{die} \de{Datei}}
\exitem{Wstaw rodzajnik określony: \de{\dots{} Programm}.}{\artn{das} \de{Programm}}
\exitem{Wstaw rodzajnik nieokreślony: \de{Das ist \dots{} Server.}}{\de{ein}}
\exitem{Wstaw rodzajnik nieokreślony: \de{Das ist \dots{} Schnittstelle.}}{\de{eine}}
\exitem{Jaki rodzaj ma \de{Anwendung} (końcówka \de{-ung})?}{żeński: \de{die Anwendung}}
\exitem{Jaki rodzaj ma \de{Mädchen} (końcówka \de{-chen})?}{nijaki: \de{das Mädchen}}
\exitem{Podaj liczbę mnogą: \de{der Stuhl}.}{\de{die Stühle}}
\exitem{Podaj liczbę mnogą: \de{die Datei}.}{\de{die Dateien}}
\exitem{Podaj liczbę mnogą: \de{das Passwort}.}{\de{die Passwörter}}
\exitem{Podaj liczbę mnogą: \de{der Computer}.}{\de{die Computer}}
\exitem{Jaki rodzaj ma \de{Netzwerk} (złożenie z \de{das Werk})?}{nijaki: \de{das Netzwerk}}
\exitem{Przetłumacz: „To jest błąd w kodzie.”}{\de{Das ist ein Fehler im Code.}}
\end{cwiczenia}

\weekend{Wybierz 15 przedmiotów z biurka i pokoju. Naklej karteczki z \emph{rodzajnikiem i słowem}
(\de{der Stuhl}, \de{die Lampe}). Opisz stanowisko pracy pięcioma zdaniami: \de{Das ist mein Arbeitsplatz. Das
ist ein Laptop.} Nagraj głos i odsłuchaj.}
